\documentclass{subfiles}
\begin{document}
    Our last section focuses on \(\np\)-completeness, specifically we focus on the \gls{tsp}.
        Although we assume the reader to be familiar with \(\np\)-problems and graph theory,
        we briefly recall what the TSP is about.

    Let \(G = (V, E)\) be a graph. We say that \(G\) has a \emph{Hamilonian cycle} if 
        there exist a close path \(P\) starting at node \(x\) and ending at \(x\) such that,
        all the vertex in \(G\) are in \(P\) and each is encountered only once.
        In other words, \(G\) has an Hamilonian cycle if we can reach each node once,
        without crossing the same one twice.

    The remainder of this section is structured as follow: 
        first we introduce the notion of approximation algorithms in the context of \(\np\)-problems,
        we then introduce an approximation algorithm for the TSP,
        and we conclude the section by showing that the generalized TSP admits no polynomial-time 
        solution, unless \(\mathcal{P} = \np\).

    \subsection{Approximation algorithms}
    \subfile{../sub/5.1 - Approximation algorithms}
\end{document}
